% ===============================
% Worksheet / Manual Template
% ===============================
\documentclass[12pt, a4paper]{article}

% ===============================
% Metadata
% ===============================
\newcommand*{\CourseCode}{MAT321}
\newcommand*{\CourseTitle}{Real Analysis II}
\newcommand*{\Chapter}{Topology of Metric Spaces}
\newcommand*{\Topics}{Open sets, Limit points and Closed sets}
\newcommand*{\DocDate}{February 15, 2026}

\include{header}

% ==========================================
% Document
% ==========================================
\begin{document}
\FrontPage
\Large


\begin{heading}
    Open and Closed Spheres
\end{heading}

\vspace{10pt}

\begin{definition}[Open Sphere]
Let $(X,d)$ be a metric space and let $a \in X$. For any $r>0$, the set
\[
S_r(a)=\{x\in X : d(x,a)<r\}
\]
is called an \emph{open sphere} (or \emph{open ball}) of radius $r$ centered at $a$.
\end{definition}

\clearpage

\begin{definition}[Closed Sphere]
Let $(X,d)$ be a metric space and let $a \in X$. For any $r>0$, the set
\[
S_r[a]=\{x\in X : d(x,a)\le r\}
\]
is called a \emph{closed sphere} of radius $r$ centered at $a$.
\end{definition}

\vspace{300pt}

\begin{problem}[Relation Between Open and Closed Spheres]
Show that
\[
S_r(a)\subseteq S_r[a]
\]
for every $a\in X$ and for every $r>0$.
\end{problem}

\clearpage

\begin{heading}
    Neighbourhood of a Point
\end{heading}

\vspace{10pt}

\begin{definition}[Neighbourhood of a Point]
Let $(X,d)$ be a metric space and let $a\in X$. A subset $N_a\subseteq X$ is called a
\emph{neighbourhood} of the point $a$ if there exists an open sphere $S_r(a)$ centered at $a$
such that
\[
S_r(a)\subseteq N_a
\]
for some $r>0$.
\end{definition}

\clearpage

\begin{problem}[Open Sphere Is a Neighbourhood of Its Points]
Show that every open sphere is a neighbourhood of each of its points.
\end{problem}

\clearpage

\begin{heading}
    Open Set
\end{heading}

\vspace{10pt}

\begin{definition}[Open Set]
A subset $G$ of a metric space $(X,d)$ is said to be \emph{open} in $X$ (with respect to the metric $d$)
if $G$ is a neighbourhood of each of its points. In other words, $G$ is open if for each $a\in G$
there exists an $r>0$ such that
\[
S_r(a)\subseteq G.
\]
\end{definition}

\clearpage

\begin{problem}[Illustrations of Open Sets]
Show that each of the following subsets is open.

\begin{enumerate}
    \item The empty set $\emptyset$ and the entire space $X$ with any metric are open sets.

    \item Every open sphere in a metric space $(X,d)$ is an open set.

    \item Let $A=\{z\in\mathbb{C} : 1<|z|<2\}$ with the usual metric. Then $A$ is an open set.

    \item Let
    \[
    S=\{(x,y)\in\mathbb{R}^2 : x^2+y^2<1\}
    \]
    with the Euclidean metric. Then $S$ is an open set.

    \item Let
    \[
    A=\{(x,y)\in\mathbb{R}^2 : y^2<x\}
    \]
    with the Euclidean metric. Then $A$ is an open set.

    \item Let
    \[
    A=\left\{f\in C[a,b] : \inf_{x\in[a,b]} f(x)>0\right\}.
    \]
    Then $A$ is open in $C[a,b]$ with respect to the supremum metric.
\end{enumerate}
\end{problem}


\clearpage

\begin{problem}[Singleton Set in $\mathbb{R}$ Is Not Open]
Show that on the real line with the usual metric, the singleton set $\{x\}$ is not open.
\end{problem}

\vspace{250pt}

\begin{problem}[Singleton Sets Under Different Metrics]
Let $\mathbb{R}$ be the set of real numbers, $d$ the usual metric and $d'$ the discrete metric on
the same set $\mathbb{R}$. Show that every singleton set $\{x\}$, $x\in\mathbb{R}$, is open in
$(\mathbb{R},d')$ but not open in $(\mathbb{R},d)$.
\end{problem}

\clearpage

\begin{problem}[Discrete Metric Space: Every Set Is Open]
Show that every subset of a discrete metric space $(X,d)$ is open.
\end{problem}

\vspace{250pt}

\begin{problem}[A Subset Open in a Subspace]
Let $(X,d)$ be a metric space where $X=(0,2)$ with the usual metric. Show that the subset
\[
A=(0,1)
\]
is an open set in $X$.
\end{problem}

\clearpage

\begin{theorem}[Basic Properties of Open Sets]
In any metric space $(X,d)$,
\begin{enumerate}
    \item the union of an arbitrary family of open sets is open;
    \item the intersection of a finite number of open sets is open.
\end{enumerate}
\end{theorem}

\clearpage

\begin{problem}[Intersection of Infinite Open Sets May Not Be Open]
Show that the intersection of an infinite number of open sets in a metric space need not be
open.
\end{problem}

\clearpage

\begin{heading}
    Limit Point
\end{heading}

\vspace{10pt}

\begin{definition}[Limit Point]
Let $(X,d)$ be a metric space and let $A\subseteq X$. A point $x\in X$ is called a
\emph{limit point} (or \emph{accumulation point}) of the set $A$ if for every $r>0$,
\[
S_r(x)\cap (A\setminus\{x\})\neq \emptyset.
\]
\end{definition}

\clearpage

\begin{problem}[Illustrations of Limit Points and Derived Sets]
Determine the set of limit points (derived set) in each of the following cases.

\begin{enumerate}
    \item Let $X=\mathbb{R}$ with the usual metric and let
    \[
    A=[0,1).
    \]
    Every point of $A$ is a limit point of $A$, and $1$ is also a limit point of $A$ although
    $1\notin A$. Hence
    \[
    A'=[0,1].
    \]

    \item Let
    \[
    A=\left\{1,\frac12,\frac13,\dots\right\}
    \]
    with the usual metric in $\mathbb{R}$. Then $0$ is the only limit point of $A$ and
    $0\notin A$. Hence
    \[
    A'=\{0\}.
    \]

    \item If $(X,d)$ is a discrete metric space, then the derived set of every subset of $X$
    is empty, that is,
    \[
    A'=\emptyset \quad \text{for all } A\subseteq X.
    \]

    \item Every real number is a limit point of the set of rational numbers $\mathbb{Q}$.
    Hence
    \[
    \mathbb{Q}'=\mathbb{R}.
    \]

    \item The set of integers $\mathbb{Z}$ has no limit point, that is,
    \[
    \mathbb{Z}'=\emptyset.
    \]

    \item Every finite subset of a metric space has no limit point.

    \item Let
    \[
    A=\{z\in\mathbb{C} : 1<|z|<2\}
    \]
    with the usual metric. Then
    \[
    A'=\{z\in\mathbb{C} : 1\le |z|\le 2\}.
    \]
\end{enumerate}
\end{problem}

\clearpage

\begin{heading}
    Closed Sets
\end{heading}

\vspace{10pt}

\begin{definition}[Closed Set]
A subset $F$ of a metric space $(X,d)$ is said to be \emph{closed} if $F$ contains all its
limit points.
\end{definition}

\clearpage

\begin{problem}[Illustrations of Closed Sets]
Verify the following statements concerning closed sets in metric spaces.

\begin{enumerate}
    \item The empty set $\emptyset$ and the whole space $X$ are closed sets in every metric
    space $(X,d)$.

    \item On the real line $\mathbb{R}$ with the usual metric, the set $\mathbb{N}$ of natural
    numbers is closed.

    \item The set $\mathbb{Q}$ of rational numbers is not closed in $\mathbb{R}$ with the usual
    metric.

    \item Every closed interval on the real line is a closed set.

    \item Every finite subset of a metric space is closed.

    \item Let
    \[
    A=\{f\in C[a,b]: f(a)=1\}.
    \]
    Then $A$ is closed in $C[a,b]$ with respect to the supremum metric.

    \item The set
    \[
    \left\{\frac{1+i}{n}: n\in\mathbb{N}\right\}
    \]
    is neither open nor closed in the complex plane with respect to the usual metric.
\end{enumerate}
\end{problem}


\clearpage

\begin{theorem}[Equivalent Definition of Closed Sets]
Let $(X,d)$ be a metric space. A subset $F\subseteq X$ is closed if and only if its complement
$X\setminus F$ is open.
\end{theorem}

\clearpage

\begin{problem}[Closed Sphere Is a Closed Set]
Show that every closed sphere in a metric space $(X,d)$ is a closed set.
\end{problem}

\clearpage

\begin{theorem}[Basic Properties of Closed Sets]
In any metric space $(X,d)$,
\begin{enumerate}
    \item the intersection of an arbitrary family of closed sets is closed;
    \item the union of a finite number of closed sets is closed.
\end{enumerate}
\end{theorem}

\clearpage

\begin{problem}[Union of Infinite Closed Sets May Not Be Closed]
Show that the union of an infinite number of closed sets in a metric space need not be
closed.
\end{problem}

\clearpage

\begin{heading}
    Closure of a Set
\end{heading}

\vspace{10pt}

\begin{definition}[Closure]
    Let $A$ be a subset of a metric space $(X,d)$. The \emph{closure} of $A$, denoted by $\overline{A}$ is the smallest closed set containing $A$, that is,
    \[
    \overline{A}=\bigcap\{F : F \text{ is closed and } A\subseteq F\};
    \]
    
\end{definition}

\clearpage

\begin{theorem}[Equivalent Definition of Closure]
Let $A$ be a subset of a metric space $(X,d)$. Then
\[
\overline{A}=A\cup A'.
\]
\end{theorem}

\clearpage

\begin{problem}[Illustrations of Closure]
Determine the closure of each of the following subsets.
\begin{enumerate}
    \item Let $X=\mathbb{R}$ with the usual metric and let
    \[
    A=(0,1)\cup(1,2).
    \]
    Then
    \[
    \overline{A}=[0,2].
    \]

    \item Let $X=\mathbb{R}$ with the usual metric and let
    \[
    A=\left\{\frac{1}{n} : n\in\mathbb{N}\right\}.
    \]
    Then
    \[
    \overline{A}=A\cup\{0\}.
    \]

    \item If $(X,d)$ is a discrete metric space and $A\subseteq X$, then
    \[
    \overline{A}=A.
    \]

    \item Let $X=\mathbb{N}$ with the usual metric and let $A=\{1,2,3\}$. Then
    \[
    \overline{A}=A.
    \]
\end{enumerate}
\end{problem}

\clearpage

\begin{theorem}[Properties of Closure]
Let $(X,d)$ be a metric space and let $A,B$ be any two subsets of $X$. Then:
\begin{enumerate}       
    \item $A$ is closed if and only if $A=\overline{A}$;
    
    \item $A\subseteq B$ implies $\overline{A}\subseteq \overline{B}$;
    
    \item $\overline{A\cup B}=\overline{A}\cup \overline{B}$;
    
    \item $\overline{A\cap B}\subseteq \overline{A}\cap \overline{B}$.
    \item $\overline{\overline{A}}=\overline{A}$.
    \item $\overline{A}=(\operatorname{ext}A)^c$.
    \item   $\overline{A}=\operatorname{int}(A^c)^c$.
\end{enumerate}
    
\end{theorem}

\clearpage




\clearpage

\EndPage
\end{document}